% GENERATED FILE. Edit canonical lessons, then run npm run generate:books.
% canonical-source-hash: fnv1a64:b876359df7f97257
% canonical-lessons: FR-C406-carrosserie, FR-C406-capot, FR-C406-portiere, FR-C406-coffre, FR-C406-banquette

\chapter{La voiture de près}
\label{ch:fr-la-voiture-de-pres}
\hlchaptermodality{406}

\begin{chapteropening}
\textbf{By the end of this chapter:} \emph{I can name the outside and inside of a car: the bodywork, the bonnet, the car door, the boot and the back seat.}

Complete the last lesson of chapter 406: I can name the outside and inside of a car: the bodywork, the bonnet, the car door, the boot and the back seat.
\end{chapteropening}

\section[la carrosserie]{la carrosserie — the bodywork (of a car)}

\label{lesson:FR-C406-carrosserie}



\begin{quote}
Before the new one: say the French for to tow, then the French for the garage owner.
\end{quote}

\subsection*{You'll want to know: la carrosserie}
\textbf{la carrosserie} — \textquotedblleft{}the bodywork (of a car)\textquotedblright{}.

\begin{quote}
La carrosserie est rayée sur le côté.
\end{quote}

\begin{quote}
\emph{The bodywork is scratched on the side.}
\end{quote}

\emph{Un carrossier} repairs bodywork.

\begin{cousinweb}
\emph{carrosserie} is built on \emph{carrosse} 'coach', which was borrowed from Italian \emph{carrozza}, from \emph{carro} and so from Latin \emph{carrus} 'wagon' (EtymDB).
\end{cousinweb}

\begin{grammarlens}[title={What you've built}]
One more word for telling, explaining and describing at B1.
\end{grammarlens}

\subsection*{Your turn}
\begin{itemize}
\raggedright
  \item \emph{Say it:} \emph{la carrosserie}
  \item \emph{Say it:} \emph{la carrosserie}, once more
  \item \emph{Say it:} a sentence of your own with \emph{la carrosserie}
  \item \emph{Recall:} say the French for to tow, then the French for the garage owner, then say \emph{la carrosserie} again
\end{itemize}

\begin{tcolorbox}[breakable,colback=teal!4,colframe=teal!35!black,title={Before you move on}]
What does la carrosserie mean? (\textquotedblleft{}The bodywork (of a car)\textquotedblright{}.) Say it once more.
\end{tcolorbox}
\section[le capot]{le capot — the bonnet (of a car)}

\label{lesson:FR-C406-capot}



\begin{quote}
Before the new one: say the French for the garage owner, then the French for the bodywork.
\end{quote}

\subsection*{You'll want to know: le capot}
\textbf{le capot} — \textquotedblleft{}the bonnet (of a car)\textquotedblright{}.

\begin{quote}
Le mécanicien ouvre le capot.
\end{quote}

\begin{quote}
\emph{The mechanic opens the bonnet.}
\end{quote}

In American English: the hood.

\begin{cousinweb}
\emph{capot} is usually traced to \emph{cape} 'cloak with a hood', from Late Latin \emph{cappa}: the engine wears a hood. EtymDB has no source line for this sense.
\end{cousinweb}

\begin{grammarlens}[title={What you've built}]
One more word for telling, explaining and describing at B1.
\end{grammarlens}

\subsection*{Your turn}
\begin{itemize}
\raggedright
  \item \emph{Say it:} \emph{le capot}
  \item \emph{Say it:} \emph{le capot}, once more
  \item \emph{Say it:} a sentence of your own with \emph{le capot}
  \item \emph{Recall:} say the French for the garage owner, then the French for the bodywork, then say \emph{le capot} again
\end{itemize}

\begin{tcolorbox}[breakable,colback=teal!4,colframe=teal!35!black,title={Before you move on}]
What does le capot mean? (\textquotedblleft{}The bonnet (of a car)\textquotedblright{}.) Say it once more.
\end{tcolorbox}
\section[la portière]{la portière — the car door}

\label{lesson:FR-C406-portiere}



\begin{quote}
Before the new one: say the French for the bodywork, then the French for the bonnet.
\end{quote}

\subsection*{You'll want to know: la portière}
\textbf{la portière} — \textquotedblleft{}the car door\textquotedblright{}.

\begin{quote}
Il a fermé la portière trop fort.
\end{quote}

\begin{quote}
\emph{He shut the car door too hard.}
\end{quote}

\emph{La porte} for a house, \emph{la portière} for a car or a train.

\begin{cousinweb}
\emph{portière} is built on \emph{porte} 'door', which is inherited from Latin \emph{porta} 'gate' (EtymDB).
\end{cousinweb}

\begin{grammarlens}[title={What you've built}]
One more word for telling, explaining and describing at B1.
\end{grammarlens}

\subsection*{Your turn}
\begin{itemize}
\raggedright
  \item \emph{Say it:} \emph{la portière}
  \item \emph{Say it:} \emph{la portière}, once more
  \item \emph{Say it:} a sentence of your own with \emph{la portière}
  \item \emph{Recall:} say the French for the bodywork, then the French for the bonnet, then say \emph{la portière} again
\end{itemize}

\begin{tcolorbox}[breakable,colback=teal!4,colframe=teal!35!black,title={Before you move on}]
What does la portière mean? (\textquotedblleft{}The car door\textquotedblright{}.) Say it once more.
\end{tcolorbox}
\section[le coffre]{le coffre — the boot (of a car); the chest}

\label{lesson:FR-C406-coffre}



\begin{quote}
Before the new one: say the French for the bonnet, then the French for the car door.
\end{quote}

\subsection*{You'll want to know: le coffre}
\textbf{le coffre} — \textquotedblleft{}the boot (of a car); the chest\textquotedblright{}.

\begin{quote}
Les valises sont dans le coffre.
\end{quote}

\begin{quote}
\emph{The suitcases are in the boot.}
\end{quote}

\emph{Un coffre} is also a big wooden box.

\begin{cousinweb}
\emph{coffre} is inherited from Latin \emph{cophinus} 'basket', which came from Greek \emph{kóphinos} (EtymDB).
\end{cousinweb}

\begin{grammarlens}[title={What you've built}]
One more word for telling, explaining and describing at B1.
\end{grammarlens}

\subsection*{Your turn}
\begin{itemize}
\raggedright
  \item \emph{Say it:} \emph{le coffre}
  \item \emph{Say it:} \emph{le coffre}, once more
  \item \emph{Say it:} a sentence of your own with \emph{le coffre}
  \item \emph{Recall:} say the French for the bonnet, then the French for the car door, then say \emph{le coffre} again
\end{itemize}

\begin{tcolorbox}[breakable,colback=teal!4,colframe=teal!35!black,title={Before you move on}]
What does le coffre mean? (\textquotedblleft{}The boot (of a car); the chest\textquotedblright{}.) Say it once more.
\end{tcolorbox}
\section[la banquette]{la banquette — the back seat; the bench seat}

\label{lesson:FR-C406-banquette}



\begin{quote}
Before the new one: say the French for the car door, then the French for the boot.
\end{quote}

\subsection*{You'll want to know: la banquette}
\textbf{la banquette} — \textquotedblleft{}the back seat; the bench seat\textquotedblright{}.

\begin{quote}
Les enfants dorment sur la banquette arrière.
\end{quote}

\begin{quote}
\emph{The children are asleep on the back seat.}
\end{quote}

\emph{La banquette arrière} — the back seat.

\begin{cousinweb}
\emph{banquette} is a small \emph{banc} 'bench', probably borrowed from Occitan \emph{banqueta}. \emph{Banc} comes from a Germanic word for 'bench' (EtymDB).
\end{cousinweb}

\begin{grammarlens}[title={What you've built}]
One more word for telling, explaining and describing at B1.
\end{grammarlens}

\subsection*{Your turn}
\begin{itemize}
\raggedright
  \item \emph{Say it:} \emph{la banquette}
  \item \emph{Say it:} \emph{la banquette}, once more
  \item \emph{Say it:} a sentence of your own with \emph{la banquette}
  \item \emph{Recall:} say the French for the car door, then the French for the boot, then say \emph{la banquette} again
\end{itemize}

\begin{tcolorbox}[breakable,colback=teal!4,colframe=teal!35!black,title={Before you move on}]
What does la banquette mean? (\textquotedblleft{}The back seat; the bench seat\textquotedblright{}.) Say it once more.
\end{tcolorbox}
